% =========================================================================
%  Figura «tipo Mathematica»: los cuatro estados elementales phi_j = 1,
%  phi_k = 0 (k != j) de un sistema de cuatro conductores en una caja a tierra.
%  Fila superior: potencial (escala azulpot de 0 a 1) con equipotenciales.
%  Fila inferior: módulo del campo (escala azulcampo, logarítmica) con las
%  mismas equipotenciales. Mapas: figuras/datos/sistema_{pot,campo}_j.png;
%  contornos: figuras/datos/sistema_eq_j.tex y sistema_formas.tex.
% =========================================================================
\pgfdeclarehorizontalshading{azulpotsh}{100bp}{rgb(0bp)=(0.961,0.973,0.988);
  rgb(20bp)=(0.863,0.906,0.961); rgb(40bp)=(0.616,0.725,0.867); rgb(60bp)=(0.227,0.431,0.647);
  rgb(80bp)=(0.122,0.306,0.537); rgb(100bp)=(0.043,0.145,0.271)}
\pgfdeclarehorizontalshading{azulcamposh}{100bp}{rgb(0bp)=(0.043,0.145,0.271);
  rgb(20bp)=(0.122,0.306,0.537); rgb(40bp)=(0.227,0.525,0.710); rgb(60bp)=(0.498,0.765,0.851);
  rgb(80bp)=(0.812,0.918,0.945); rgb(100bp)=(1,1,1)}
\newcommand{\panelSistema}[3]{% #1 = j (1..4), #2 = pot / campo, #3 = conductor activo (A..D)
  \node[anchor=south west, inner sep=0] at (-5,-4.5)
    {\includegraphics[width=3.7cm, height=3.33cm]{figuras/datos/sistema_#2_#1.png}};
  \begin{scope}
    \clip (-5,-4.5) rectangle (5,4.5);
    \input{figuras/datos/sistema_eq_#1}
  \end{scope}
  \foreach \f in {\formaA,\formaB,\formaC,\formaD}{\fill[white, draw=azulprof, line width=0.4pt] \f;}
  \fill[azulprof, draw=azulprof, line width=0.4pt] \csname forma#3\endcsname;
  \draw[azulprof, line width=0.5pt] (-5,-4.5) rectangle (5,4.5);}

\begin{figure}[t]
\centering
\begin{tikzpicture}[x=0.37cm, y=0.37cm, font=\footnotesize]
  % Contornos de los cuatro conductores
% Generado por scripts/generar_datos.py (no editar a mano).
\def\formaA{(-4.250,-2.950) -- (-3.050,-3.950) -- (0.220,-0.180) -- cycle}
\def\formaB{(-2.750,1.250) -- (-0.650,1.250) -- (-0.650,2.850) -- (-2.750,2.850) -- cycle}
\def\formaC{(1.330,0.350) -- (1.330,0.363) -- (1.329,0.375) -- (1.329,0.388) -- (1.327,0.400) -- (1.326,0.413) -- (1.324,0.425) -- (1.322,0.437) -- (1.320,0.450) -- (1.317,0.462) -- (1.314,0.474) -- (1.310,0.486) -- (1.307,0.498) -- (1.302,0.510) -- (1.298,0.522) -- (1.293,0.534) -- (1.289,0.545) -- (1.283,0.557) -- (1.278,0.568) -- (1.272,0.579) -- (1.266,0.590) -- (1.259,0.601) -- (1.253,0.611) -- (1.246,0.622) -- (1.238,0.632) -- (1.231,0.642) -- (1.223,0.652) -- (1.215,0.662) -- (1.207,0.671) -- (1.198,0.680) -- (1.189,0.689) -- (1.180,0.698) -- (1.171,0.707) -- (1.162,0.715) -- (1.152,0.723) -- (1.142,0.731) -- (1.132,0.738) -- (1.122,0.746) -- (1.111,0.753) -- (1.101,0.759) -- (1.090,0.766) -- (1.079,0.772) -- (1.068,0.778) -- (1.057,0.783) -- (1.045,0.789) -- (1.034,0.793) -- (1.022,0.798) -- (1.010,0.802) -- (0.998,0.807) -- (0.986,0.810) -- (0.974,0.814) -- (0.962,0.817) -- (0.950,0.820) -- (0.937,0.822) -- (0.925,0.824) -- (0.913,0.826) -- (0.900,0.827) -- (0.888,0.829) -- (0.875,0.829) -- (0.863,0.830) -- (0.850,0.830) -- (0.837,0.830) -- (0.825,0.829) -- (0.812,0.829) -- (0.800,0.827) -- (0.787,0.826) -- (0.775,0.824) -- (0.763,0.822) -- (0.750,0.820) -- (0.738,0.817) -- (0.726,0.814) -- (0.714,0.810) -- (0.702,0.807) -- (0.690,0.802) -- (0.678,0.798) -- (0.666,0.793) -- (0.655,0.789) -- (0.643,0.783) -- (0.632,0.778) -- (0.621,0.772) -- (0.610,0.766) -- (0.599,0.759) -- (0.589,0.753) -- (0.578,0.746) -- (0.568,0.738) -- (0.558,0.731) -- (0.548,0.723) -- (0.538,0.715) -- (0.529,0.707) -- (0.520,0.698) -- (0.511,0.689) -- (0.502,0.680) -- (0.493,0.671) -- (0.485,0.662) -- (0.477,0.652) -- (0.469,0.642) -- (0.462,0.632) -- (0.454,0.622) -- (0.447,0.611) -- (0.441,0.601) -- (0.434,0.590) -- (0.428,0.579) -- (0.422,0.568) -- (0.417,0.557) -- (0.411,0.545) -- (0.407,0.534) -- (0.402,0.522) -- (0.398,0.510) -- (0.393,0.498) -- (0.390,0.486) -- (0.386,0.474) -- (0.383,0.462) -- (0.380,0.450) -- (0.378,0.437) -- (0.376,0.425) -- (0.374,0.413) -- (0.373,0.400) -- (0.371,0.388) -- (0.371,0.375) -- (0.370,0.363) -- (0.370,0.350) -- (0.370,0.337) -- (0.371,0.325) -- (0.371,0.312) -- (0.373,0.300) -- (0.374,0.287) -- (0.376,0.275) -- (0.378,0.263) -- (0.380,0.250) -- (0.383,0.238) -- (0.386,0.226) -- (0.390,0.214) -- (0.393,0.202) -- (0.398,0.190) -- (0.402,0.178) -- (0.407,0.166) -- (0.411,0.155) -- (0.417,0.143) -- (0.422,0.132) -- (0.428,0.121) -- (0.434,0.110) -- (0.441,0.099) -- (0.447,0.089) -- (0.454,0.078) -- (0.462,0.068) -- (0.469,0.058) -- (0.477,0.048) -- (0.485,0.038) -- (0.493,0.029) -- (0.502,0.020) -- (0.511,0.011) -- (0.520,0.002) -- (0.529,-0.007) -- (0.538,-0.015) -- (0.548,-0.023) -- (0.558,-0.031) -- (0.568,-0.038) -- (0.578,-0.046) -- (0.589,-0.053) -- (0.599,-0.059) -- (0.610,-0.066) -- (0.621,-0.072) -- (0.632,-0.078) -- (0.643,-0.083) -- (0.655,-0.089) -- (0.666,-0.093) -- (0.678,-0.098) -- (0.690,-0.102) -- (0.702,-0.107) -- (0.714,-0.110) -- (0.726,-0.114) -- (0.738,-0.117) -- (0.750,-0.120) -- (0.763,-0.122) -- (0.775,-0.124) -- (0.787,-0.126) -- (0.800,-0.127) -- (0.812,-0.129) -- (0.825,-0.129) -- (0.837,-0.130) -- (0.850,-0.130) -- (0.863,-0.130) -- (0.875,-0.129) -- (0.888,-0.129) -- (0.900,-0.127) -- (0.913,-0.126) -- (0.925,-0.124) -- (0.937,-0.122) -- (0.950,-0.120) -- (0.962,-0.117) -- (0.974,-0.114) -- (0.986,-0.110) -- (0.998,-0.107) -- (1.010,-0.102) -- (1.022,-0.098) -- (1.034,-0.093) -- (1.045,-0.089) -- (1.057,-0.083) -- (1.068,-0.078) -- (1.079,-0.072) -- (1.090,-0.066) -- (1.101,-0.059) -- (1.111,-0.053) -- (1.122,-0.046) -- (1.132,-0.038) -- (1.142,-0.031) -- (1.152,-0.023) -- (1.162,-0.015) -- (1.171,-0.007) -- (1.180,0.002) -- (1.189,0.011) -- (1.198,0.020) -- (1.207,0.029) -- (1.215,0.038) -- (1.223,0.048) -- (1.231,0.058) -- (1.238,0.068) -- (1.246,0.078) -- (1.253,0.089) -- (1.259,0.099) -- (1.266,0.110) -- (1.272,0.121) -- (1.278,0.132) -- (1.283,0.143) -- (1.289,0.155) -- (1.293,0.166) -- (1.298,0.178) -- (1.302,0.190) -- (1.307,0.202) -- (1.310,0.214) -- (1.314,0.226) -- (1.317,0.238) -- (1.320,0.250) -- (1.322,0.263) -- (1.324,0.275) -- (1.326,0.287) -- (1.327,0.300) -- (1.329,0.312) -- (1.329,0.325) -- (1.330,0.337) -- cycle}
\def\formaD{(1.050,1.320) -- (3.550,1.320) -- (3.550,1.680) -- (1.050,1.680) -- cycle}

  \tikzset{numEQ/.style={azulprof!80, line width=0.25pt}}
  \foreach \j/\c/\l in {1/A/a,2/B/b,3/C/c,4/D/d}{
    \begin{scope}[xshift={(\j-1)*3.95cm}]
      \panelSistema{\j}{pot}{\c}
      \node[anchor=south] at (0,4.55) {(\l) $\phi_{\mathrm{\c}}=1$};
    \end{scope}
    \begin{scope}[xshift={(\j-1)*3.95cm}, yshift=-3.75cm]
      \tikzset{numEQ/.style={white!85!azulprof, line width=0.25pt}}
      \panelSistema{\j}{campo}{\c}
    \end{scope}}
  % barras de color
  \begin{scope}[yshift=-5.9cm]
    \shade[shading=azulpotsh] (-5,-0.2) rectangle (14.2,0.35);
    \draw[azulclaro, line width=0.3pt] (-5,-0.2) rectangle (14.2,0.35);
    \node[anchor=north] at (-5,-0.25) {$0$};
    \node[anchor=north] at (14.2,-0.25) {$1$};
    \node[anchor=north] at (4.6,-0.25) {potencial $\phi$};
    \shade[shading=azulcamposh] (17.8,-0.2) rectangle (37.03,0.35);
    \draw[azulclaro, line width=0.3pt] (17.8,-0.2) rectangle (37.03,0.35);
    \node[anchor=north] at (17.8,-0.25) {débil};
    \node[anchor=north] at (37.03,-0.25) {intenso};
    \node[anchor=north] at (27.4,-0.25) {$\norm{\vE}$ (escala logarítmica)};
  \end{scope}
\end{tikzpicture}
\caption{Estados elementales que definen los coeficientes de capacidad e influencia de un sistema de cuatro conductores (A~cuña, B~rectángulo, C~disco
y D~barra) dentro de una caja a tierra, calculados numéricamente. En cada columna, el conductor $j$ (relleno oscuro) está a potencial unidad y los demás a
tierra: la carga de cada conductor en ese estado es la columna $j$ de $\Cm$. Arriba, el potencial; abajo, el módulo del campo, que es máximo junto a las
aristas y en los huecos estrechos entre conductores próximos.}
\label{fig:sistema-cij}
\end{figure}
